\documentclass[../../main.tex]{subfiles}

\begin{document}

\section{度量拓扑（续）}\label{section:ch02-度量拓扑续}

在这一节中，我们讨论度量拓扑与前面引进的一些概念之间的关系。

正如我们所期望的那样，度量空间的子空间还是度量空间。如果 $A$ 是拓扑空间 $X$ 的一个子空间，$d$ 是 $X$ 上的一个度量，那么 $d$ 在 $A \times A$ 上的限制便是诱导 $A$ 的拓扑的一个度量。证明留给读者完成。

我们不打算就序拓扑作进一步的讨论。易见有些序拓扑是可度量化的（如 $\mathbb{Z}_{+}$ 与 $\mathbb{R}$），有些则不然。

每一个度量拓扑都满足 Hausdorff 公理。若 $x$ 和 $y$ 是度量空间 $(X, d)$ 中不同的两点，令 $\varepsilon = \frac{1}{2}d(x, y)$，则由三角不等式可得 $B_{d}(x, \varepsilon)$ 与 $B_{d}(y, \varepsilon)$ 无交。

积拓扑的特殊情形已经研究过，证明了 $\mathbb{R}^{n}$ 与 $\mathbb{R}^{\omega}$ 是可度量化的。一般说来，可度量化空间的可数积是可度量化的。其证明仿照 $\mathbb{R}^{\omega}$ 可度量化的证明进行，我们留作习题。

这里将对连续函数给予特别的关注。本节余下的部分将研究这个问题。

当我们考虑度量空间上的连续函数时，本书所进行的讨论很接近于微积分和分析学，有两件事要做：

第一，希望将我们所熟知的连续性的“$\varepsilon$-$\delta$ 定义”推广到一般度量空间上，并且对连续性的“收敛序列定义”也做同样的工作。

第二，除了在第\cref{section:ch02-连续函数}节中讨论过的那些构造连续函数的方法之外，还希望研究两种构造连续函数的方法。一个是取连续实值函数的和、差、积、商。另一个是取一致收敛的连续函数序列的极限。

\begin{theorem}[$\varepsilon$-$\delta$ 连续性]\label{theorem:ch02-epsdelta连续}
设 $f: X \rightarrow Y$，$X$ 和 $Y$ 是分别具有度量 $d_{X}$ 和 $d_{Y}$ 的可度量化空间。那么 $f$ 的连续性等价于：对给定的 $x \in X$ 和 $\varepsilon > 0$，存在 $\delta > 0$，使得
\begin{align*}
d_{X}(x, y) < \delta \Longrightarrow d_{Y}(f(x), f(y)) < \varepsilon.
\end{align*}
\end{theorem}

\begin{proof}
设 $f$ 连续。给定 $x$ 和 $\varepsilon$。这时，集合
\begin{align*}
f^{-1}(B(f(x), \varepsilon))
\end{align*}
是 $X$ 的含有 $x$ 的开集，从而它包含某一个以 $x$ 为中心的 $\delta$-球 $B(x, \delta)$。若 $y$ 在这个 $\delta$-球中，则 $f(y)$ 在以 $f(x)$ 为中心的 $\varepsilon$-球中。

反之，设 $\varepsilon$-$\delta$ 条件满足，令 $V$ 是 $Y$ 的一个开集，我们证明 $f^{-1}(V)$ 在 $X$ 中是开的。设 $x$ 是 $f^{-1}(V)$ 的一个点。因为 $f(x) \in V$，所以存在一个以 $f(x)$ 为中心的 $\varepsilon$-球 $B(f(x), \varepsilon)$ 包含于开集 $V$。根据 $\varepsilon$-$\delta$ 条件，存在以 $x$ 为中心的 $\delta$-球 $B(x, \delta)$，使得 $f(B(x, \delta)) \subset B(f(x), \varepsilon)$。于是 $B(x, \delta)$ 是包含于 $f^{-1}(V)$ 的点 $x$ 的一个邻域，因此 $f^{-1}(V)$ 是一个开集。
\end{proof}

再来看连续性的收敛序列定义。我们首先讨论收敛序列与集合的闭包间的关系。根据分析学中的经验，人们确信：若点 $x$ 是 $X$ 的子集 $A$ 的闭包中的点，那么存在一个 $A$ 的点的一个序列收敛到 $x$。一般而言，这是不对的，但在度量空间中这个结论成立。

\begin{lemma}[序列引理]\label{lemma:ch02-序列引理}
设 $X$ 是一个拓扑空间，$A \subset X$。若 $A$ 中有一个收敛于 $x$ 的序列，则 $x \in \overline{A}$。若 $X$ 为可度量化空间，则逆命题也成立。
\end{lemma}

\begin{proof}
设 $x_{n} \rightarrow x$，其中 $x_{n} \in A$。这时 $x$ 的任意邻域都包含 $A$ 的一个点，根据\cref{theorem:ch02-闭包邻域}，$x \in \overline{A}$。

反之，设 $X$ 是可度量化的，$x \in \overline{A}$。令 $d$ 为诱导出 $X$ 的拓扑的一个度量。对于任意正整数 $n$，取以 $x$ 为中心 $1/n$ 为半径的邻域 $B_{d}(x, \frac{1}{n})$。在它与 $A$ 的交中选取点 $x_{n}$。我们证明序列 $x_{n}$ 收敛于 $x$。这是因为任何包含 $x$ 的开集 $U$，必定包含一个以 $x$ 为中心的 $\varepsilon$-球 $B_{d}(x, \varepsilon)$，如果取足够大的 $N$，使得 $\frac{1}{N} < \varepsilon$，那么对于所有 $i \geqslant N$ 有 $x_{i}$ 属于 $U$。
\end{proof}

\begin{theorem}[序列连续性]\label{theorem:ch02-序列连续}
设 $f: X \rightarrow Y$。若 $X$ 为可度量化的空间，则 $f$ 连续的充分必要条件是对于 $X$ 中每一个收敛序列 $x_{n} \rightarrow x$，序列 $f(x_{n})$ 收敛于 $f(x)$。
\end{theorem}

\begin{proof}
设 $f$ 连续。给定 $x_{n} \rightarrow x$，我们证明 $f(x_{n}) \rightarrow f(x)$。令 $V$ 是 $f(x)$ 的一个邻域，则 $f^{-1}(V)$ 是 $x$ 的一个邻域。于是存在 $N$，使得对于 $n \geqslant N$ 有 $x_{n} \in f^{-1}(V)$。因此对 $n \geqslant N$，$f(x_{n}) \in V$。

我们证明充分性。设收敛序列的条件满足。若 $A$ 是 $X$ 的一个子集，下面证明 $f(\overline{A}) \subset \overline{f(A)}$。如果 $x \in \overline{A}$，则根据上一个引理可见：存在 $A$ 中点的序列 $x_{n}$ 收敛于 $x$。根据假设，序列 $f(x_{n})$ 收敛于 $f(x)$。因为 $f(x_{n}) \in f(A)$，根据上一个引理可见 $f(x) \in \overline{f(A)}$。（注意，并不要求 $Y$ 是可度量化的。）这就证明了 $f(\overline{A}) \subset \overline{f(A)}$。
\end{proof}

附带指出，在证明引理\cref{lemma:ch02-序列引理}和定理\cref{theorem:ch02-序列连续}时，并没有用到 $X$ 是可度量化空间这么强的条件，而只要求在 $x$ 处有球 $B_{d}(x, 1/n)$ 的一个可数族。这启发我们引入以下定义。

\begin{definition}[可数基与第一可数性公理]\label{definition:ch02-第一可数性}
空间 $X$ 称为在 $x$ 处有可数基（countable basis），如果存在 $x$ 的可数邻域族 $\{U_{n}\}_{n \in \mathbb{Z}_{+}}$，使得 $x$ 的任意邻域 $U$ 至少包含一个 $U_{n}$。如果空间 $X$ 在每一点处有一个可数基，则称 $X$ 满足第一可数性公理（first countability axiom）。
\end{definition}

若 $X$ 在 $x$ 处有一个可数基，则引理\cref{lemma:ch02-序列引理}的证明仍可完成。只要用集合
\begin{align*}
B_{n} = U_{1} \cap U_{2} \cap \cdots \cap U_{n}
\end{align*}
代替球 $B_{d}(x, \frac{1}{n})$ 便可以了。定理\cref{theorem:ch02-序列连续}的证明也可以不加改变地完成。

一个可度量化空间总满足第一可数性公理，我们将会看到其逆不真。为了证明有关空间的一些定理，就像 Hausdorff 公理那样，有时必须在拓扑空间上加上第一可数性公理。第 4 章将详细地讨论相关内容。

现在研究构造连续函数的其他方法，这需要用到以下引理：

\begin{lemma}[算术运算的连续性]\label{lemma:ch02-算术连续}
加法、减法和乘法运算是从 $\mathbb{R} \times \mathbb{R}$ 到 $\mathbb{R}$ 中的连续函数，除法运算是从 $\mathbb{R} \times (\mathbb{R} - \{0\})$ 到 $\mathbb{R}$ 中的连续函数。
\end{lemma}

读者也许会发现这个引理在前面已经证明过了，就是标准的“$\varepsilon$-$\delta$ 论证”。你也可以在下面的习题 12 中找到证明的要点。因此，不难写出证明的细节。

\begin{theorem}[连续实值函数的运算]\label{theorem:ch02-实值函数运算}
设 $X$ 是一个拓扑空间，$f, g: X \rightarrow \mathbb{R}$ 连续，则 $f + g$、$f - g$、$f \cdot g$ 都是连续的。若对于所有 $x$，$g(x) \neq 0$，则 $f/g$ 也是连续的。
\end{theorem}

\begin{proof}
根据\cref{theorem:ch02-积空间映射}，用式子
\begin{align*}
h(x) = f(x) \times g(x)
\end{align*}
定义的映射 $h: X \rightarrow \mathbb{R} \times \mathbb{R}$ 是连续的。函数 $f + g$ 是 $h$ 和加法运算
\begin{align*}
+: \mathbb{R} \times \mathbb{R} \longrightarrow \mathbb{R}
\end{align*}
的复合，因而 $f + g$ 连续。对于 $f - g$、$f \cdot g$ 及 $f/g$ 证明是类似的。
\end{proof}

最后，我们来讨论一致收敛的概念。

\begin{definition}[一致收敛]\label{definition:ch02-一致收敛}
设 $f_{n}: X \rightarrow Y$ 是一个从集合 $X$ 到度量空间 $Y$ 的函数序列，$d$ 是 $Y$ 中的度量。我们称序列 $(f_{n})$ 一致收敛（converges uniformly）于函数 $f: X \rightarrow Y$，如果对于任意给定的 $\varepsilon > 0$，存在整数 $N$，对于 $n > N$ 以及任意 $x \in X$ 有
\begin{align*}
d(f_{n}(x), f(x)) < \varepsilon.
\end{align*}
\end{definition}

收敛的一致性不仅依赖于 $Y$ 的拓扑，也依赖于 $Y$ 上的度量。我们有以下关于一致收敛序列的定理：

\begin{theorem}[一致极限定理]\label{theorem:ch02-一致极限定理}
设 $f_{n}: X \rightarrow Y$ 是从拓扑空间 $X$ 到度量空间 $Y$ 的连续函数的一个序列。若 $(f_{n})$ 一致收敛于 $f$，则 $f$ 连续。
\end{theorem}

\begin{proof}
设 $V$ 是 $Y$ 的一个开集，$x_{0}$ 为 $f^{-1}(V)$ 的一个点。我们要找出 $x_{0}$ 的一个邻域 $U$，使得 $f(U) \subset V$。

令 $y_{0} = f(x_{0})$。先取 $\varepsilon$，使得 $\varepsilon$-球 $B(y_{0}, \varepsilon)$ 包含在 $V$ 中，然后应用一致收敛性选取 $N$，使得对于所有 $n \geqslant N$ 以及任意 $x \in X$ 有
\begin{align*}
d(f_{n}(x), f(x)) < \varepsilon/3.
\end{align*}
最后再应用 $f_{N}$ 的连续性，选取 $x_{0}$ 的一个邻域 $U$，使得 $f_{N}$ 将 $U$ 映射到 $Y$ 的以 $f_{N}(x_{0})$ 为中心的 $\varepsilon/3$ 球内。

下面证明 $f$ 将 $U$ 映射到 $B(y_{0}, \varepsilon)$ 中，从而映射到 $V$ 中，这便是我们所要证明的。为此，我们注意当 $x \in U$ 时，有
\begin{align*}
d(f(x), f_{N}(x)) < \varepsilon/3 \quad \text{（根据 } N \text{ 的选取）}, \\
d(f_{N}(x), f_{N}(x_{0})) < \varepsilon/3 \quad \text{（根据 } U \text{ 的选取）}, \\
d(f_{N}(x_{0}), f(x_{0})) < \varepsilon/3 \quad \text{（根据 } N \text{ 的选取）}.
\end{align*}
把这三个式子相加，并且应用三角不等式，就得到所要证明的
\begin{align*}
d(f(x), f(x_{0})) < \varepsilon.
\end{align*}
\end{proof}

注意一致收敛概念与上节所讲的一致度量的定义有关系。例如，考虑由全体函数 $f: X \rightarrow \mathbb{R}$ 的集合连同其上一致度量 $\bar{\rho}$ 所形成的空间 $\mathbb{R}^{X}$。不难看出，函数序列 $f_{n}: X \rightarrow \mathbb{R}$ 一致收敛于 $f$ 的充分必要条件是当把 $f_{n}$ 看成度量空间 $(\mathbb{R}^{X}, \bar{\rho})$ 的元素时，序列 $(f_{n})$ 收敛于 $f$。有关的证明留作习题。

作为本节的结束，我们给出一些不可度量化空间的例子。

\par
\begin{figure}[htbp]
\centering
\includegraphics[width=0.5\textwidth]{ch02-fig21-1}
\caption{函数 $f_{n}(x) = \frac{1}{n^{3}[x - (1/n)]^{2} + 1}$}\label{figure:ch02-度量拓扑续-21.1}
\end{figure}

\begin{example}\label{example:ch02-Romega箱不可度量}
$\mathbb{R}^{\omega}$ 在箱拓扑下是不可度量化的。

我们将证明序列引理对 $\mathbb{R}^{\omega}$ 不成立。设 $A$ 为 $\mathbb{R}^{\omega}$ 中所有坐标都是正数的点所组成的那个子集，即
\begin{align*}
A = \{(x_{1}, x_{2}, \cdots) \mid \text{对于所有 } i \in \mathbb{Z}_{+}, x_{i} > 0\}.
\end{align*}
设 $\boldsymbol{0}$ 是 $\mathbb{R}^{\omega}$ 的“原点”，即零点 $(0, 0, \cdots)$，它的每一个坐标都是 0。在箱拓扑中，$\boldsymbol{0} \in \overline{A}$，这是因为，如果
\begin{align*}
B = (a_{1}, b_{1}) \times (a_{2}, b_{2}) \times \cdots
\end{align*}
是含有 $\boldsymbol{0}$ 的任意一个基元素，则 $B$ 与 $A$ 相交，比如点
\begin{align*}
\left(\frac{1}{2}b_{1}, \frac{1}{2}b_{2}, \cdots\right)
\end{align*}
就属于 $B \cap A$。

但是，我们可以断定在 $A$ 中没有收敛于 $\boldsymbol{0}$ 的点的序列。对于 $A$ 中的序列 $(\boldsymbol{a}_{n})$，
\begin{align*}
\boldsymbol{a}_{n} = (x_{1n}, x_{2n}, \cdots, x_{in}, \cdots),
\end{align*}
其中每一个坐标 $x_{in} > 0$，这就可以构造出 $\mathbb{R}$ 的箱拓扑的基元素 $B'$，
\begin{align*}
B' = (-x_{11}, x_{11}) \times (-x_{22}, x_{22}) \times \cdots.
\end{align*}
$B'$ 含有原点 $\boldsymbol{0}$，却不包含序列 $(\boldsymbol{a}_{n})$ 的元素：这是因为 $\boldsymbol{a}_{n}$ 的第 $n$ 个坐标 $x_{nn}$ 不属于区间 $(-x_{nn}, x_{nn})$。所以点 $\boldsymbol{a}_{n}$ 不可能属于 $B'$。因此在箱拓扑下，序列 $\boldsymbol{a}_{n}$ 不收敛于 $\boldsymbol{0}$。
\end{example}

\begin{example}\label{example:ch02-不可数积不可度量}
不可数个 $\mathbb{R}$ 的积空间是不可度量化的。

设 $J$ 为不可数指标集。我们来证明：$\mathbb{R}^{J}$ 在积拓扑下不满足序列引理。

设 $A$ 为 $\mathbb{R}^{J}$ 的子集，其元素 $(x_{\alpha})$ 满足：除了有限个 $\alpha$ 之外，对其余所有的 $\alpha$ 有 $x_{\alpha} = 1$。设 $\boldsymbol{0}$ 为 $\mathbb{R}^{J}$ 的“原点”，它的所有坐标都是 0。

可以论断 $\boldsymbol{0}$ 属于 $A$ 的闭包。事实上，令 $\prod U_{\alpha}$ 为含有 $\boldsymbol{0}$ 的一个基元素，那么仅对有限个 $\alpha$，比如说 $\alpha = \alpha_{1}, \cdots, \alpha_{n}$，有 $U_{\alpha} \neq \mathbb{R}$。设 $(x_{\alpha})$ 为 $A$ 中一点，对于 $\alpha = \alpha_{1}, \cdots, \alpha_{n}$，取 $x_{\alpha} = 0$。对其余所有的 $\alpha$，取 $x_{\alpha} = 1$。那么 $(x_{\alpha}) \in A \cap \prod U_{\alpha}$，这便证明了 $\boldsymbol{0} \in \overline{A}$。

但是 $A$ 中不存在收敛于 $\boldsymbol{0}$ 的点的序列。我们对此证明如下：设 $(\boldsymbol{a}_{n})$ 为 $A$ 的点的一个序列。给定 $n$，设 $J_{n}$ 表示指标集 $J$ 的子集，其元素 $\alpha \in J_{n}$ 是使 $\boldsymbol{a}_{n}$ 的第 $\alpha$ 个坐标异于 1 的指标。集合 $J_{n}$ 的并是有限集的可数并，因而是可数的。但是 $J$ 不可数，那么 $J$ 中必有一个指标，比如说 $\beta$，使得它不在每一个 $J_{n}$ 中。这就是说：对于每一个点 $\boldsymbol{a}_{n}$，都有它的第 $\beta$ 个坐标等于 1。

令 $U_{\beta}$ 为 $\mathbb{R}$ 中的开区间 $(-1, 1)$，$U$ 为 $\mathbb{R}^{J}$ 的一个开集 $\pi_{\beta}^{-1}(U_{\beta})$。那么 $U$ 是 $\boldsymbol{0}$ 的一个邻域，并且它不包含任何点 $\boldsymbol{a}_{n}$。所以序列 $(\boldsymbol{a}_{n})$ 不可能收敛于 $\boldsymbol{0}$。
\end{example}

\subsection*{习题}

\begin{exercise}\label{exercise:ch02-度量拓扑续-1}
设 $A \subset X$。若 $d$ 是 $X$ 的拓扑的一个度量，证明 $d \mid A \times A$ 是 $A$ 的子空间拓扑的一个度量。
\end{exercise}

\begin{exercise}\label{exercise:ch02-度量拓扑续-2}
设 $X$ 和 $Y$ 是分别具有度量 $d_{X}$ 和 $d_{Y}$ 的度量空间。对于 $X$ 中任意两点 $x_{1}$ 和 $x_{2}$，$f: X \rightarrow Y$ 满足条件
\begin{align*}
d_{Y}(f(x_{1}), f(x_{2})) = d_{X}(x_{1}, x_{2}).
\end{align*}
证明：$f$ 是一个嵌入。$f$ 称为 $X$ 到 $Y$ 中的一个等距嵌入（isometric imbedding）。
\end{exercise}

\begin{exercise}\label{exercise:ch02-度量拓扑续-3}
假设对于 $n \in \mathbb{Z}_{+}$，$X_{n}$ 是以 $d_{n}$ 为度量的度量空间。

(a) 证明
\begin{align*}
\rho(x, y) = \max \{d_{1}(x_{1}, y_{1}), \cdots, d_{n}(x_{n}, y_{n})\}
\end{align*}
是积空间 $X_{1} \times \cdots \times X_{n}$ 上的一个度量。

(b) 设 $\bar{d}_{i} = \min \{d_{i}, 1\}$。证明
\begin{align*}
D(x, y) = \sup \{\bar{d}_{i}(x_{i}, y_{i})/i\}
\end{align*}
是积空间 $\prod X_{i}$ 上的一个度量。
\end{exercise}

\begin{exercise}\label{exercise:ch02-度量拓扑续-4}
证明 $\mathbb{R}_{l}$ 和有序矩形都满足第一可数公理。（当然，这并不意味着它们是可度量化的。）
\end{exercise}

\begin{exercise}\label{exercise:ch02-度量拓扑续-5}
在空间 $\mathbb{R}$ 中，设 $x_{n} \rightarrow x$ 且 $y_{n} \rightarrow y$，则
\begin{align*}
x_{n} + y_{n} \rightarrow x + y, \\
x_{n} - y_{n} \rightarrow x - y, \\
x_{n} y_{n} \rightarrow xy,
\end{align*}
并且如果每一个 $y_{n} \neq 0$ 和 $y \neq 0$ 则有
\begin{align*}
x_{n}/y_{n} \rightarrow x/y.
\end{align*}
[提示：应用引理\cref{lemma:ch02-算术连续}。根据第\cref{section:ch02-积拓扑}节中的习题可见：如果 $x_{n} \rightarrow x$ 和 $y_{n} \rightarrow y$，则有 $x_{n} \times y_{n} \rightarrow x \times y$。]
\end{exercise}

\begin{exercise}\label{exercise:ch02-度量拓扑续-6}
$f_{n}: [0, 1] \rightarrow \mathbb{R}$ 定义为 $f_{n}(x) = x^{n}$。证明：对于每一个 $x \in [0, 1]$ 序列 $(f_{n}(x))$ 收敛，但是序列 $(f_{n})$ 不致使收敛。
\end{exercise}

\begin{exercise}\label{exercise:ch02-度量拓扑续-7}
设 $X$ 是一个集合，$f_{n}: X \rightarrow \mathbb{R}$ 是函数的一个序列。$\bar{\rho}$ 是空间 $\mathbb{R}^{X}$ 上的一致度量。证明：序列 $(f_{n})$ 一致收敛于函数 $f: X \rightarrow \mathbb{R}$ 当且仅当把 $f_{n}$ 看成度量空间 $(\mathbb{R}^{X}, \bar{\rho})$ 的元素时，序列 $(f_{n})$ 收敛于 $f$。
\end{exercise}

\begin{exercise}\label{exercise:ch02-度量拓扑续-8}
设 $X$ 是一个拓扑空间，$Y$ 是一个度量空间。$f_{n}: X \rightarrow Y$ 是连续函数的一个序列，$x_{n}$ 是 $X$ 中收敛于 $x$ 的一个点的序列。证明：若序列 $(f_{n})$ 一致收敛于 $f$，则 $(f_{n}(x_{n}))$ 收敛于 $f(x)$。
\end{exercise}

\begin{exercise}\label{exercise:ch02-度量拓扑续-9}
设函数 $f_{n}: \mathbb{R} \rightarrow \mathbb{R}$ 定义为
\begin{align*}
f_{n}(x) = \frac{1}{n^{3}[x - (1/n)]^{2} + 1}.
\end{align*}
参见\cref{figure:ch02-度量拓扑续-21.1}。设 $f: \mathbb{R} \rightarrow \mathbb{R}$ 为零函数。

(a) 证明：对于任何 $x \in \mathbb{R}$，$f_{n}(x) \rightarrow f(x)$。

(b) 证明：$f_{n}$ 不一致收敛于 $f$。（这说明定理\cref{theorem:ch02-一致极限定理}的逆不成立，即，虽然 $f_{n}$ 不一致收敛于 $f$，其极限函数 $f$ 也可能连续。）
\end{exercise}

\begin{exercise}\label{exercise:ch02-度量拓扑续-10}
应用连续性的闭集式定义（\cref{theorem:ch02-连续等价}），证明下列集合是 $\mathbb{R}^{2}$ 的闭子集
\begin{align*}
A = \{x \times y \mid xy = 1\}, \\
S^{1} = \{x \times y \mid x^{2} + y^{2} = 1\}, \\
B^{2} = \{x \times y \mid x^{2} + y^{2} \leqslant 1\}.
\end{align*}
集合 $B^{2}$ 称为 $\mathbb{R}^{2}$ 的（闭）单位球（unit ball）。
\end{exercise}

\begin{exercise}\label{exercise:ch02-度量拓扑续-11}
证明无穷级数的下列基本性质。

(a) 若 $(s_{n})$ 是一个有界实数列，并且对于任意 $n$ 有 $s_{n} \leqslant s_{n+1}$，则 $(s_{n})$ 收敛。

(b) 设 $(a_{n})$ 是一个实数列，
\begin{align*}
s_{n} = \sum_{i=1}^{n} a_{i}.
\end{align*}
若 $s_{n} \rightarrow s$，则称无穷级数（infinite series）
\begin{align*}
\sum_{i=1}^{\infty} a_{i}
\end{align*}
收敛于 $s$。证明：若 $\sum a_{i}$ 收敛于 $s$，$\sum b_{i}$ 收敛于 $t$，则 $\sum (ca_{i} + b_{i})$ 收敛于 $cs + t$。

(c) 证明无穷级数的比较判别法（comparison test）：若对于每一个 $i$ 有 $|a_{i}| \leqslant b_{i}$，并且级数 $\sum b_{i}$ 收敛，则级数 $\sum a_{i}$ 收敛。[提示：证明级数 $\sum |a_{i}|$ 与 $\sum c_{i}$ 收敛，其中 $c_{i} = |a_{i}| + a_{i}$。]

(d) 给定一个函数序列 $f_{n}: X \rightarrow \mathbb{R}$，令
\begin{align*}
s_{n}(x) = \sum_{i=1}^{n} f_{i}(x).
\end{align*}
证明一致收敛性的 Weierstrass $M$-判别法（Weierstrass $M$-test）：若对于所有的 $x \in X$ 和 $i$ 有 $|f_{i}(x)| \leqslant M_{i}$，并且级数 $\sum M_{i}$ 收敛，则序列 $(s_{n})$ 一致收敛于一个函数 $s$。[提示：令 $r_{n} = \sum_{i=n+1}^{\infty} M_{i}$。证明：若 $k > n$，则 $|s_{k}(x) - s_{n}(x)| \leqslant r_{n}$。因此，$|s(x) - s_{n}(x)| \leqslant r_{n}$。]
\end{exercise}

\begin{exercise}\label{exercise:ch02-度量拓扑续-12}
证明 $\mathbb{R}$ 中的代数运算的连续性如下：应用 $\mathbb{R}$ 的度量 $d(a, b) = |a - b|$ 和 $\mathbb{R}^{2}$ 的度量
\begin{align*}
\rho((x, y), (x_{0}, y_{0})) = \max \{|x - x_{0}|, |y - y_{0}|\}.
\end{align*}
(a) 证明加法运算连续。[提示：任意给定 $\varepsilon$，令 $\delta = \varepsilon/2$，并且注意
\begin{align*}
d(x + y, x_{0} + y_{0}) \leqslant |x - x_{0}| + |y - y_{0}|.]
\end{align*}
(b) 证明乘法运算连续。[提示：任意给定 $(x_{0}, y_{0})$ 及 $0 < \varepsilon < 1$，令
\begin{align*}
3\delta = \varepsilon/(|x_{0}| + |y_{0}| + 1),
\end{align*}
并且注意
\begin{align*}
d(xy, x_{0}y_{0}) \leqslant |x_{0}|\, |y - y_{0}| + |y_{0}|\, |x - x_{0}| + |x - x_{0}|\, |y - y_{0}|.]
\end{align*}
(c) 证明取倒数运算是从 $\mathbb{R} - \{0\}$ 到 $\mathbb{R}$ 的连续映射。[提示：证明 $(a, b)$ 的原像是开集。就 $a$ 和 $b$ 为正、为负和为零考虑五种情况。]

(d) 证明减法运算和除法运算是连续的。
\end{exercise}

\end{document}